\documentclass[11pt]{article}
\usepackage[margin=1in]{geometry}
\usepackage{amsmath,amssymb,amsthm,booktabs}
\usepackage[no-math]{fontspec}
\usepackage{fancyvrb,microtype}
\usepackage[colorlinks=true,allcolors=blue]{hyperref}

\setmonofont{DejaVuSansMono.ttf}[Scale=MatchLowercase]
\fvset{fontsize=\small,xleftmargin=1em}

\theoremstyle{definition}
\newtheorem{theorem}{Theorem}[section]
\newtheorem{lemma}[theorem]{Lemma}

\newcommand{\R}{\mathbb{R}}
\newcommand{\Z}{\mathbb{Z}}
\newcommand{\Zcal}{\mathcal{Z}}
\newcommand{\0}{\mathbf{0}}
\newcommand{\1}{\mathbf{1}}
\newcommand{\x}{\mathbf{x}}
\newcommand{\y}{\mathbf{y}}
\newcommand{\z}{\mathbf{z}}
\newcommand{\s}{\mathbf{s}}
\newcommand{\e}{\mathbf{e}}
\newcommand{\w}{\mathbf{w}}
\renewcommand{\v}{\mathbf{v}}
\renewcommand{\b}{\mathbf{b}}
\DeclareMathOperator{\conv}{conv}
\DeclareMathOperator{\vol}{vol}

\newcommand{\notetitle}{A counterexample to Oertel\texorpdfstring{\textquotesingle}{'}s conjecture on Mixed-integer volume}
\title{\vspace{-2.5em}\notetitle\\[0.4em]
\large The sharp finite-dimensional form}
\author{}
\date{}
\hypersetup{pdftitle={\notetitle}}

\begin{document}
\maketitle
\vspace{-3.5em}

\begin{abstract}
We construct a rational polytope $C\subset\R^{12}\times\R^5$ of positive mixed-integer volume and a
number $r<(5/6)^5/2^{12}$ such that every point of $\Z^{12}\times\R^5$ lies in a proper closed
halfspace whose intersection with $C$ has at most the fraction $r$ of that volume. The polytope has
$4852$ nonempty integer fibers. Two selector coordinates and a sign change on four of the other ten
integer coordinates produce three families of $252$ light fibers. Each heavy fiber paired with a
light fiber consists of a simplex and a slab of another simplex sharing a facet. The volume
comparisons use exact rational arithmetic. Lean~4 verifies the uniform depth bound.
\end{abstract}

\section{The result}\label{sec:result}

Conjecture~17 of Basu and Oertel~\cite{basu-oertel} asserts that every convex body in
$\R^n\times\R^d$ of positive mixed-integer volume has a point in $\Z^n\times\R^d$ of normalized
halfspace depth at least $2^{-n}(d/(d+1))^d$. We disprove this bound at $(n,d)=(12,5)$, where it
equals $(5/6)^5/2^{12}$. Cheng and Basu~\cite{cheng-basu} use the name Oertel's conjecture in their
Conjecture~2.3 for the lower bound $1/(e2^n)$.

Let $\vol$ denote Lebesgue measure on $\R^5$. The \emph{mixed-integer volume}, or \emph{mass}, of a set
$D\subseteq\R^{12}\times\R^5$ with measurable slices is
$\mu(D)=\sum_{\z\in\Z^{12}}\vol\{\x\in\R^5:(\z,\x)\in D\}$.

\begin{theorem}\label{thm:main}
There exist a rational polytope $C\subset\R^{12}\times\R^5$ with nonempty interior, $\0\in C$ and
$0<\mu(C)<\infty$, and a real number $r$ with $0<r<(5/6)^5/2^{12}$, such that the following holds. For
every $(\z,\x)\in\Z^{12}\times\R^5$ there are $b\in\R$ and a nonzero linear functional $F$ on
$\R^{12}\times\R^5$ with $F(\z,\x)\ge b$ and $\mu(C\cap\{F\ge b\})\le r\,\mu(C)$.
\end{theorem}

The \emph{depth} of a mixed-integer point is the infimum of $\mu(C\cap H)/\mu(C)$ over proper closed
halfspaces $H$ containing it. The theorem bounds the supremum of these depths by $r$. In units of
the fiber over the origin, the construction gives a containing halfspace of mass at most $(5h/6)^5$
at every mixed-integer point, while the total mass $W$ exceeds $4096h^5$. This gives $r=(5h/6)^5/W$.

\section{Continuous fibers}\label{sec:fibers}

Throughout, we fix the rational numbers
\begin{gather*}
a=\frac{663}{1000},\qquad \delta=\frac{664201}{10^6},\qquad \varepsilon=\frac{13}{20000},\qquad h=\frac{1111946415}{10^9},\qquad \eta=\Bigl(\frac{h}{1+2\delta}\Bigr)^5,\\
(t_0,\ldots,t_5)=10^{-9}(166735371,\,166735294,\,166324902,\,166734811,\,166734811,\,166734811).
\end{gather*}
They satisfy $\sum_it_i=1$, $1/2<a<1$, $0<\delta<1$, $\varepsilon<1/1000$ and $1/7<t_i<1/5$. Also
$2h<1+2\delta<3$, which gives $\eta<1/32<a/(1+a)$. The only other numerical facts we use are the strict
comparisons at the end of Section~\ref{sec:cuts}.

Let $B=\{\x\in\R^6:x_i\ge0,\ \sum_ix_i=1\}$ be the standard simplex with vertices
$\v_1,\ldots,\v_6$. We identify the hyperplane $\sum_ix_i=1$ with $\R^5$ by dropping $x_6$ and
subtracting $1/6$ from the other coordinates. The centroid of $B$ is then the origin of $\R^5$, scalar
multiples of subsets are dilations about the centroid, and $x_1,\ldots,x_6$ are affine functions on
$\R^5$ with sum $1$. From now on $\vol$ and $\mu$ are measured in units of the Lebesgue volume of $B$,
which does not change the ratios in Theorem~\ref{thm:main}.

For $j=1,2,3$, with these three indices taken cyclically, put $f_j=x_j-ax_{j+1}-(1-a)x_{j+2}$. Each
$f_j$ vanishes at the centroid and equals $1$ at $\v_j$, so it is a nonzero linear functional on $\R^5$.
Moreover $|f_j|\le1$ on $B$ and $f_1+f_2+f_3=0$. The cyclic permutation of $\v_1,\v_2,\v_3$ is a
volume-preserving linear symmetry of $B$ that carries $f_j$ to $f_{j+1}$, so we describe the sets below
for $j=1$ and obtain the others by rotation.

Each $f_j$ is a homogeneous linear form in $(x_1,x_2,x_3)$. Non-strict inequalities $f_j\ge0$ or
$f_j\le0$ therefore cut out a join of their trace on the active triangle $\conv(\v_1,\v_2,\v_3)$
with $\conv(\v_4,\v_5,\v_6)$. Indeed, normalizing the first three coordinates by their sum when it is
positive preserves the inequalities. Points with that sum zero belong to the latter triangle. The simplex spanned by three points of the active triangle and $\v_4,\v_5,\v_6$
has volume equal to the absolute determinant of the $3\times3$ matrix of barycentric coordinates of these
points, because the affine map carrying the vertices of $B$ to the vertices of this simplex acts on
barycentric coordinates as the direct sum of this matrix and the identity. We call this the determinant
formula. On the active triangle,
$f_1$ is positive at $\v_1$, negative at $\v_2$ and $\v_3$, and zero at $(a\v_1+\v_2)/(1+a)$ and
$((1-a)\v_1+\v_3)/(2-a)$. So $B\cap\{f_1\ge0\}$ is a simplex, and the determinant formula gives its volume
$\beta=1/[(1+a)(2-a)]$.

Let $T_j=B\cap\{f_{j+1}\le0,\ f_{j+2}\le0\}$, so that $f_j\ge0$ on $T_j$ and $f_i\le0$ on $T_j$ for
$i\ne j$. Among the vertices of the active triangle, $f_2$ is positive only at $\v_2$ and $f_3$ only at
$\v_3$. The lines $f_2=0$ and $f_3=0$ meet at the centroid $\bar\v=(\v_1+\v_2+\v_3)/3$ of the triangle.
They meet the edges from $\v_1$ to $\v_2$ and from $\v_1$ to $\v_3$ at $\w_2=(\v_1+(1-a)\v_2)/(2-a)$ and
$\w_3=(\v_1+a\v_3)/(1+a)$, respectively. The edge from $\v_2$ to $\v_3$ has $f_1<0$ throughout, so
it is disjoint from $T_1$. The trace of $T_1$ is the quadrilateral with vertices
$\v_1,\w_2,\bar\v,\w_3$, whose diagonal from $\v_1$ to $\bar\v$ lies on the line $x_2=x_3$. Consequently
$T_1=U\cup P$, where
\[
U=\conv(\v_1,\bar\v,\w_3,\v_4,\v_5,\v_6),\qquad P=\conv(\v_1,\bar\v,\w_2,\v_4,\v_5,\v_6).
\]
These simplices lie on opposite sides of the hyperplane $x_2=x_3$ and meet in a common facet, which is a
null set. By the determinant formula, $U$ has volume $\theta=a/[3(1+a)]$ and $P$ has volume
$\phi=(1-a)/[3(2-a)]$.

To truncate $T_1$, let $\nu_0,\ldots,\nu_5$ be the barycentric coordinates of $U$ in the listed vertex
order and put $q_i=\nu_i(\w_2)$. Comparing active coordinates gives $q_1=3(1-a)/(2-a)$,
$q_2=-(1-a^2)/[a(2-a)]$, $q_0=1-q_1-q_2$ and $q_3=q_4=q_5=0$. Since $1/2<a<1$, we have $0<q_1<1$, $q_2<0$
and $0<q_1+q_2=q_1(2a-1)/(3a)<1$, so $0<q_0<1$. Now put $K_1=T_1\cap\{\nu_2\ge\varepsilon q_2\}$. On $U$
we have $\nu_2\ge0>\varepsilon q_2$. On $P$, let $p$ be the barycentric coordinate at $\w_2$ and, for
$i\ne2$, let $p_i$ be the barycentric coordinate at the vertex of $U$ with index $i$. Then
$\nu_i=p_i+q_ip$ for $i\ne2$ and $\nu_2=q_2p$, since both sides are affine and agree at the vertices of
$P$. Therefore $K_1$ is the union of $U$ and $P\cap\{p\le\varepsilon\}$. The set $P\cap\{p\ge\varepsilon\}$
is a homothetic copy of $P$ with ratio $1-\varepsilon$, so $\vol(K_1)=V:=\theta+\phi[1-(1-\varepsilon)^5]$.
Rotation gives $K_2,K_3$ and their coordinates $\nu_i$.

Finally, let $N_j=\conv\bigl(\eta\v_j+(1-\eta)\v_{j+1},\v_{j+1},\v_{j+2},\v_4,\v_5,\v_6\bigr)\subseteq B$.
At these vertices $f_j$ takes the values $\eta-a(1-\eta)$, $-a$, $a-1$, $0$, $0$, $0$, which are
nonpositive because $\eta<a/(1+a)$. Hence $f_j\le0$ on $N_j$, and $\vol(N_j)=\eta$ by the determinant
formula.

\section{Integer bases and the polytope}\label{sec:bases}

Write $\z=(\y,\s)\in\R^{10}\times\R^2$, with ordinary coordinates $\y$ and selector coordinates $\s$, and
put $u=\sum_iy_i$. The linear forms that organize the construction are
\[
\tau(\z)=(u-s_2,\ u-s_1,\ u-s_1-s_2),\qquad g(\z)=\delta\bigl(\tau_1(\z)+\tau_2(\z)+\tau_3(\z)\bigr)=\delta(3u-2s_1-2s_2).
\]
Let $c_i=-1$ for $i\le4$ and $c_i=1$ for $i\ge5$. For $J\subseteq\{1,\ldots,10\}$, let $\1_J$ be the
indicator vector of $J$, and let $\b_J$ have entries $c_i$ on $J$ and $0$ off $J$. With $\e_1,\e_2$ the unit
vectors of $\R^2$, the candidates are $(\1_J,\0)$ for $J\ne\emptyset$ and $(\b_J,\s)$ for
$\s\in\{\e_1,\e_2,\e_1+\e_2\}$, the latter with the following values.
\begin{center}
\begin{tabular}{cccccc}
\toprule
Type $j$ & $\s$ & $\tau(\b_J,\s)$ & Sum of entries & Paired & Near \\
\midrule
$1$ & $\e_1$ & $(u,\ u-1,\ u-1)$ & $3u-2$ & $u=1$ & $u=0$ \\
$2$ & $\e_2$ & $(u-1,\ u,\ u-1)$ & $3u-2$ & $u=1$ & $u=0$ \\
$3$ & $\e_1+\e_2$ & $(u-1,\ u-1,\ u-2)$ & $3u-4$ & $u=1$ & $u=2$ \\
\bottomrule
\end{tabular}
\end{center}
The sum of the entries of $\tau$ is not divisible by $3$ in this table and equals $3|J|$ at $(\1_J,\0)$, so
it is nonzero on every candidate. A \emph{heavy base} is a candidate multiplied by the sign of this sum.
After this orientation, the paired value of $u$ in type $j$ gives $\tau_j=1$ and zeros elsewhere, the near
value gives $\tau_j=0$ and ones elsewhere, and every other value of $u$ gives three positive entries, as do
the bases $(\1_J,\0)$. We call these heavy bases \emph{paired} of type $j$, \emph{near} of type $j$ and
\emph{regular}, respectively. For each paired base $\z$ of type $j$, the point $-\z$ is a \emph{light} base
of type $j$. The set $\Zcal$ of \emph{bases} consists of $\0$ and all heavy and light bases.

Modulo $2$, an oriented $(\b_J,\s)$ reduces to $(\1_J,\s)$, and the three selectors have distinct nonzero
residues. Hence the heavy bases represent each nonzero class of $(\Z/2\Z)^{12}$ exactly once, and there are
$4095$ of them. Since $u(\b_J)=|J\mathbin{\triangle}\{1,2,3,4\}|-4$, toggling membership in the first four
positions shows that $u=1$ has $\binom{10}{5}=252$ solutions in each type. So there are $756$ paired bases,
$756$ light bases and $3339$ other heavy bases. The same count at $u=0$ and $u=2$ gives
$\binom{10}{4}=\binom{10}{6}=210$ near bases per type, leaving $2709$ regular bases. A light base
$-\z$ has its partner's residue and differs from it, so the $4852$ points of $\Zcal$ are distinct.
We assign to them the following fibers.
\begin{center}
\begin{tabular}{llcllc}
\toprule
Base $\z$ & Number & $g(\z)$ & $\tau(\z)$ & Fiber $A_{\z}$ & $\vol(A_{\z})$ \\
\midrule
$\0$ & $1$ & $0$ & $\0$ & $B$ & $1$ \\
paired, type $j$ & $252$ & $\delta$ & $\tau_j=1$, others $0$ & $(1+\delta)K_j$ & $(1+\delta)^5V$ \\
light, type $j$ & $252$ & $-\delta$ & $\tau_j=-1$, others $0$ & $(1-\delta)T_j$ & $(1-\delta)^5(\theta+\phi)$ \\
near, type $j$ & $210$ & $2\delta$ & $\tau_j=0$, others $1$ & $(1+2\delta)N_j$ & $h^5$ \\
regular & $2709$ & $\ge3\delta$ & all entries $\ge1$ & $hB$ & $h^5$ \\
\bottomrule
\end{tabular}
\end{center}
The volume of a near fiber is $(1+2\delta)^5\eta=h^5$. Every fiber lies in $(1+g(\z))B$, and the radius
$1+g(\z)$ is positive because $\delta<1$. For regular bases this containment uses $h<1+2\delta<1+g(\z)$ and
$\0\in B$. Let $C=\conv\bigl(\bigcup_{\z\in\Zcal}\{\z\}\times A_{\z}\bigr)$. The light fibers of one type
have total volume $L=252(\theta+\phi)(1-\delta)^5$, and all fibers have total volume
$W=3339h^5+756(1+\delta)^5V+1+3L$.

\section{The integer fibers of \texorpdfstring{$C$}{C}}\label{sec:census}

We first show that a nonzero integer point of $\conv(\Zcal)$ has only the trivial representation.

\begin{lemma}\label{lem:saturation}
If $\z\in\Z^{12}\setminus\{\0\}$ is a convex combination $\sum_a\lambda_a\z^a$ of points $\z^a\in\Zcal$
with all $\lambda_a>0$, then $\z^a=\z$ for every $a$.
\end{lemma}

\begin{proof}
All entries of all bases lie in $\{-1,0,1\}$, so $\z=(\y,\s)$ has the same property. If a coordinate of
$\z$ equals $\pm1$, or if the contributors $\z^a=(\y^a,\s^a)$ take only two consecutive integer values in
that coordinate, then integrality forces all contributors to agree in it. We call this saturation.

Suppose first that $\s\ne\0$, say $s_k=\pm1$. Saturation gives $s_k^a=s_k$ for all $a$. Every nonzero
selector of a base is $\pm\e_1$, $\pm\e_2$ or $\pm(\e_1+\e_2)$, so the other selector coordinate of each
contributor lies in $\{0,s_k\}$, and saturation gives $\s^a=\s$. Let $\omega=\pm1$ be the sign of the
nonzero entries of $\s$. By construction, a base with this selector has ordinary part $\omega\b_J$, so
$y_i^a\in\{0,\omega c_i\}$ for each $i$, and saturation gives $\y^a=\y$.

Now suppose that $\s=\0$. Then $y_i=\pm1$ for some $i$, and saturation gives $y_i^a=y_i$ for all $a$. A
contributor with a nonzero selector of sign $\omega$ has $y_i=\omega c_i$, so
$c_iy_i(s_1^a+s_2^a)=\omega(s_1^a+s_2^a)>0$, while a contributor with selector $\0$ gives $0$. The weighted
sum of these numbers is $c_iy_i(s_1+s_2)=0$, so all contributors have selector $\0$. Their ordinary parts
lie in $\{0,1\}^{10}$, and saturation gives $\y^a=\y$.
\end{proof}

Consequently $\Z^{12}\cap\conv(\Zcal)=\Zcal$, and no nonzero base lies in the convex hull of the other
bases. For a point $\w$ outside a compact convex set, let $\w_0$ be a nearest point in the set.
Minimizing the squared distance to $\w$ on segments from $\w_0$ gives
$\langle\w-\w_0,\w'-\w_0\rangle\le0$ for every point $\w'$ of the set. The inner product with
$\w-\w_0$ is thus strictly larger at $\w$ than on the set. In particular, each nonzero
$\z\in\Zcal$ has a linear functional $\ell_{\z}$ on
$\R^{12}$ with $\ell_{\z}(\z')<\ell_{\z}(\z)$ for all $\z'\in\Zcal\setminus\{\z\}$.

\begin{lemma}\label{lem:fibers}
The set $C$ is a rational polytope with nonempty interior and $\0\in C$. For $\z\in\Z^{12}$, the slice
$\{\x:(\z,\x)\in C\}$ equals $A_{\z}$ if $\z\in\Zcal$ and is empty otherwise. Consequently $\mu(C)=W$, and
every closed halfspace $H$ satisfies $\mu(C\cap H)=\sum_{\z\in\Zcal}\vol\bigl(A_{\z}\cap\{\x:(\z,\x)\in H\}\bigr)$.
\end{lemma}

\begin{proof}
The fibers are polytopes with rational vertices, because $T_j$ and $K_j$ are cut from $B$ by rational
inequalities, $N_j$ has rational vertices, and the dilation factors are rational. So $C$ is the convex hull
of finitely many rational points. It contains $\{\0\}\times B$, hence the origin. The set $B$ spans the
continuous directions. The regular bases $(\1_{\{i\}},\0)$ and the near bases $(\0,-\e_1)$ and $(\0,-\e_2)$
span the integer directions, so $C$ has nonempty interior.

Let $(\z,\x)\in C$ with $\z$ integer, and write $(\z,\x)=\sum_a\lambda_a(\z^a,\x^a)$ with $\lambda_a>0$,
$\sum_a\lambda_a=1$, $\z^a\in\Zcal$ and $\x^a\in A_{\z^a}$. If $\z\ne\0$, then Lemma~\ref{lem:saturation}
gives $\z^a=\z$ for all $a$, so $\z\in\Zcal$ and $\x\in A_{\z}$ by convexity of $A_{\z}$. If $\z=\0$, then
$\x^a/(1+g(\z^a))\in B$ for all $a$. Since $g$ is linear, the positive weights $\lambda_a(1+g(\z^a))$ sum to
$1+g(\z)=1$. Thus $\x$ is a convex combination of the points $\x^a/(1+g(\z^a))$ of $B$ with these weights,
which gives $\x\in B=A_{\0}$. The reverse inclusions are immediate, and the formulas for $\mu$ follow by
summing over $\Zcal$.
\end{proof}

\section{Halfspaces of small mass}\label{sec:cuts}

\begin{proof}[Proof of Theorem~\ref{thm:main}]
By Lemma~\ref{lem:fibers}, a point $(\z,\x)\in\Z^{12}\times\R^5$ either lies outside $C$ or has
$\z\in\Zcal$ and $\x\in A_{\z}$. All inequalities defining the covers below are nonstrict, so points on the
boundaries of fibers need no separate treatment.

\noindent\emph{Exterior points.} If $(\z,\x)\notin C$, nearest-point separation gives a nonzero $F$ with
$F(\z,\x)>F$ on $C$. With $b=F(\z,\x)$, the halfspace $\{F\ge b\}$ contains the point and misses $C$, so it
has mass $0$.

\noindent\emph{The origin and the light fibers.} For $j=1,2,3$, let
$H_j=\{(\z,\x):f_j(\x)\ge3\tau_j(\z)\}$, a proper closed halfspace because $f_j\ne0$. We read its trace on
each fiber from the table in Section~\ref{sec:bases}. Over the origin the trace is $B\cap\{f_j\ge0\}$, of
volume $\beta$. If $\tau_j(\z)=-1$, the condition $f_j\ge-3$ holds on the whole light fiber of type $j$. If
$\z\ne\0$ and $\tau_j(\z)=0$, then $\z$ is a paired or light base of a type $k\ne j$, or a near base of type
$j$. Its fiber lies in $\{f_j\le0\}$, so the trace lies in the hyperplane $\{f_j=0\}$, which is null because
$f_j$ is nonconstant. If $\tau_j(\z)\ge1$, the condition requires $f_j\ge3$, which fails because every heavy
fiber lies in $(1+2\delta)B$, where $|f_j|\le1+2\delta<3$. Hence $\mu(C\cap H_j)=\beta+L$. Since
$f_1+f_2+f_3=0$, each $\x\in B$ has $f_j(\x)\ge0$ for some $j$, so $(\0,\x)\in H_j$. A point of a light fiber
of type $j$ lies in $H_j$.

\noindent\emph{Isolated cuts.} Let $\z\ne\0$ be a base, let $\psi$ be a nonconstant affine function on
$\R^5$, and let $\kappa\in\R$. For $M>0$ put $F(\z',\x')=\psi(\x')-\psi(\0)+M\ell_{\z}(\z')$ and
$b=\kappa-\psi(\0)+M\ell_{\z}(\z)$. Then $F$ is a nonzero linear functional, and $F(\z,\x')\ge b$ if and only
if $\psi(\x')\ge\kappa$. Since $\Zcal$ is finite, the fibers are compact and $\ell_{\z}(\z')<\ell_{\z}(\z)$
for $\z'\ne\z$, we have $F<b$ on every fiber over $\z'\ne\z$ once $M$ is large. By Lemma~\ref{lem:fibers},
the isolated cut $\{F\ge b\}$ then has mass $\vol(A_{\z}\cap\{\psi\ge\kappa\})$.

\noindent\emph{Regular and near fibers.} Such a fiber is a simplex of volume $h^5$. Some barycentric
coordinate of $\x$ in it is at least $1/6$. That coordinate is affine and nonconstant, and the cap of the
simplex where it is at least $1/6$ is a homothetic copy with ratio $5/6$. So the isolated cut has mass
$(5h/6)^5$.

\noindent\emph{Paired fibers.} Let $\z$ be a paired base of type $j$ and let $\x\in(1+\delta)K_j$. By rotation
we compute with $j=1$. Since $\sum_i\nu_i=\sum_it_i=1$ on all of $\R^5$, there is an index $i$ with
$\nu_i(\x/(1+\delta))\ge t_i$. The function $\x'\mapsto\nu_i(\x'/(1+\delta))$ is affine and nonconstant, so
the isolated cut with $\kappa=t_i$ contains $(\z,\x)$ and has mass $(1+\delta)^5V_i$, where
$V_i=\vol(K_1\cap\{\nu_i\ge t_i\})$.

To compute $V_i$, we split $K_1$ into $U$ and $P\cap\{p\le\varepsilon\}$, which overlap in a null set. In $U$
the cap $\{\nu_i\ge t_i\}$ is a homothetic copy of $U$ with ratio $1-t_i$, of volume $\theta(1-t_i)^5$. On $P$ with $i=2$ we have
$\nu_2=q_2p\le0<t_2$, so this part is empty. Let $i\ne2$. For $0\le p\le\varepsilon$, the bounds on $q_i$ and
$t_i$ give $t_i-q_ip>0$, $1-t_i-(1-q_i)p>0$ and $1-q_i>0$. In the coordinates $p$ and $p_k$ with $k\ne i,2$,
the simplex $P$ is a standard simplex of Lebesgue volume $1/120$, and its volume in our units is $\phi$.
Substituting $p_i=1-p-\sum_{k\ne i,2}p_k$ turns the condition $\nu_i=p_i+q_ip\ge t_i$ into
$\sum_{k\ne i,2}p_k\le1-t_i-(1-q_i)p$, and it forces $p_i\ge t_i-q_ip>0$. So for each $p$ the four remaining
coordinates range over a simplex of Lebesgue volume $[1-t_i-(1-q_i)p]^4/24$, and the contribution of $P$ is
\[
5\phi\int_0^\varepsilon\bigl[1-t_i-(1-q_i)p\bigr]^4\,dp=\phi\,\frac{(1-t_i)^5-\bigl(1-t_i-(1-q_i)\varepsilon\bigr)^5}{1-q_i}.
\]
Thus $V_i$ is $\theta(1-t_i)^5$ plus this quantity for $i\ne2$, and $V_2=\theta(1-t_2)^5$.

\noindent\emph{Exact comparisons.} The cuts above reduce the theorem to comparisons of explicit rational
numbers. In lowest terms, $\theta=221/1663$, $\phi=337/4011$, $\beta=1000000/2223431$,
$q_0=776569/886431$, $q_1=1011/1337$, $q_2=-560431/886431$ and $5h/6=74129761/80000000$. Exact evaluation of
the formulas above gives the following floors, and the checker in Appendix~\ref{sec:lean} repeats these
comparisons with exact fractions.
\begin{center}
\begin{tabular}{lc}
\toprule
Margin & $\lfloor10^{12}\cdot\text{margin}\rfloor$ \\
\midrule
$(5h/6)^5-\beta-L$ & $1262439$ \\
$(5h/6)^5-(1+\delta)^5V_i$ for $i=0,\ldots,5$ & $7542,\ 6288,\ 1211,\ 7869,\ 7869,\ 7869$ \\
$W-4096h^5$ & $45459473$ \\
\bottomrule
\end{tabular}
\end{center}
All floors are positive, so $\beta+L<(5h/6)^5$, $(1+\delta)^5V_i<(5h/6)^5$ for $i=0,\ldots,5$, and
$4096h^5<W$. Therefore every point of $\Z^{12}\times\R^5$ lies in a proper closed halfspace $H$ with
$\mu(C\cap H)\le(5h/6)^5$. Put $r=(5/6)^5h^5/W$. Since $\mu(C)=W$, we get $\mu(C\cap H)\le r\,\mu(C)$.
Since $W>4096h^5$, we also get $0<r<(5/6)^5/2^{12}$. Numerically, $W=6962.7250497\ldots$,
$10^5r=9.811464096\ldots$ and $10^5(5/6)^5/2^{12}=9.811464160\ldots$. Lemma~\ref{lem:fibers} gives the
remaining properties of $C$.
\end{proof}

\begin{thebibliography}{9}
\bibitem{basu-oertel}
Amitabh Basu and Timm Oertel.
\emph{\href{https://www.ams.jhu.edu/~abasu9/papers/centerpoints.pdf}{Centerpoints:
A link between optimization and convex geometry}}.
Manuscript version with Conjecture~17.

\bibitem{cheng-basu}
Hongyu Cheng and Amitabh Basu.
\emph{\href{https://hongyucheng.net/files/linear-threshold-oertel.pdf}{Linear Threshold
for Oertel's Conjecture on the Mixed-Integer Volume}}.
Conjecture~2.3.
\end{thebibliography}

\clearpage
\appendix
\section{Lean statement and exact replay}\label{sec:lean}

The exact checker and Lean project are tracked in
\href{https://github.com/Hongyu-Cheng/weak-oertel-conjecture/pull/14}{PR~\#14}.
The verification uses the following pins.
\begin{center}
\begin{tabular}{ll}
\toprule
Verification source & \texttt{d368306f28ce808d2b2b773cc025745d835538bf} \\
Lean toolchain & \texttt{leanprover/lean4:v4.34.1} \\
mathlib revision & \texttt{d13f23b723b8a846827a245b89c10fc7d3f11612} \\
\bottomrule
\end{tabular}
\end{center}
In the namespace \texttt{SharpOC.Selector17}, the final declaration reads as follows.
\begin{Verbatim}
theorem sharp_oc_counterexample :
    ∃ C : Set Ambient, IsCompact C ∧ Convex ℝ C ∧ (interior C).Nonempty ∧
      ((0 : IntegerSpace), (0 : E)) ∈ C ∧
      0 < mixedIntegerMass C ∧ mixedIntegerMass C ≠ ⊤ ∧
      ∃ r : ℝ, 0 < r ∧ r < (((5 : ℝ) / 6)^5 / (2 : ℝ)^12) ∧
        ∀ (z : Coordinate → ℤ) (x : E),
          ∃ F : Ambient →L[ℝ] ℝ, ∃ b : ℝ, F ≠ 0 ∧ b ≤ F ((fun j => (z j : ℝ)), x) ∧
            (mixedIntegerMass (C ∩ {p | b ≤ F p})).toReal /
              (mixedIntegerMass C).toReal ≤ r
\end{Verbatim}
Here \texttt{Coordinate} is \texttt{Fin 10 ⊕ Fin 2}, \texttt{IntegerSpace} is \texttt{Coordinate → ℝ},
\texttt{E} is \texttt{Fin 5 → ℝ} and \texttt{Ambient} is \texttt{IntegerSpace × E}. For a set \texttt{D}, the
value \texttt{mixedIntegerMass D} is the extended nonnegative sum, over all \texttt{w : Coordinate → ℤ}, of the
Lebesgue volume of \texttt{\{x | ((fun k => (w k : ℝ)), x) ∈ D\}}, which is $\mu(D)$ in the Lebesgue units of
Section~\ref{sec:result}. Finiteness of the mass makes both \texttt{toReal} values faithful.
Lemma~\ref{lem:fibers} establishes rationality of the polytope. With Python 3, Git and elan installed,
run the following from the repository root. Lake reads its dependency pins from
\texttt{lean/lake-manifest.json}.
\begin{Verbatim}
set -eu
python3 -B research/sharp-dimension-17/check.py
cd lean
lake exe cache get
lake build
lake env lean CheckAxioms.lean
lake env lean --stdin <<'LEAN'
import SharpOC.Selector17Counterexample

#check SharpOC.Selector17.integer_projection_census
#check SharpOC.Selector17.ambient_dimension
#check SharpOC.Selector17.sharp_oc_counterexample
#print axioms SharpOC.Selector17.sharp_oc_counterexample
LEAN
\end{Verbatim}
The checker uses exact rational arithmetic to enumerate the $4852$ listed bases and check the strict
comparisons. The geometric proof and Lean establish that the list exhausts the integer projection.
The Lean build passed, and the dimension lemma and final theorem depend only on \texttt{propext},
\texttt{Classical.choice} and \texttt{Quot.sound}.
\end{document}
